\chapter{Optimal Execution: Almgren--Chriss}
\label{ch:almgren_chriss}
\chaptermark{Almgren--Chriss optimal execution}

You need to sell $1{,}000{,}000$ shares of a stock by the end of the
day. Selling them all at once walks the order book and moves the
market against you by tens of basis points. Trailing them out
one-by-one over the full session exposes you to all-day price risk,
where a random adverse drift of half a percent on tens of millions
of dollars notional wipes out a week of desk P\&L. Between the two
extremes sits a schedule (how many shares to sell at each minute)
that balances rushing against dawdling. Under linear impact
and Gaussian price noise, ``when should you sell each share'' has
a clean closed-form answer, and this chapter derives it from scratch.

The answer is the \textbf{Almgren--Chriss (A--C) optimal execution
trajectory} \citeAC{}, the benchmark every sell-side execution desk
measures algorithms against. It \emph{assumes linear market impact},
directly contradicting the empirical square-root law of
\cref{ch:market_impact}; reconciling the two is part of this
chapter's job. The objective resembles the inventory-averse market
maker of \cref{ch:market_making} (the dual problem: earn spread
subject to inventory risk vs.\ pay impact subject to price risk on
unsold inventory), but here the closed form is in elementary
functions and we derive it by calculus of variations on finite sums.

\begin{backgroundbox}[Prerequisites]
From earlier in the book:
\begin{itemize}
  \item \Cref{ch:market_impact}: the empirical square-root law
  $I(Q) \approx Y\volat\sqrt{Q/V}$; temporary (reverting) vs.\
  permanent (information-revealing) impact; linear impact
  over-predicts cost by a factor
  $\sqrt{Q_{\text{exec}}/Q_{\text{calib}}}$ when extrapolated out
  of calibration range. This chapter assumes linear impact for
  tractability; the contradiction is flagged in
  \cref{sec:ac_vs_sqrt}.
  \item \Cref{ch:market_making}: the Avellaneda--Stoikov
  inventory-averse market maker. A market maker \emph{earns}
  spread subject to price risk on inventory; the A--C
  liquidator \emph{pays} impact subject to price risk on
  not-yet-liquidated inventory. Same mean-variance objective
  and HJB machinery, opposite ends of the client--dealer
  relationship. Warning: A--S uses $\riskav$ for the
  market-maker's risk aversion; here $\riskav$ is the
  permanent-impact coefficient (pitfallbox below).
  \item \Cref{ch:mean_reversion_ou}: mean-variance thinking
  (fix an expected cost, minimise variance, or vice versa). A--C
  applies it to a schedule rather than a scalar.
\end{itemize}
From outside the book:
\begin{itemize}
  \item Elementary \textbf{Lagrange multipliers} / discrete
  calculus of variations: differentiating a quadratic in a
  trajectory $\{x_k\}$ under a linear terminal constraint.
  \item \textbf{Linear recurrences}: the recurrence
  $x_{k-1} - 2\cosh(\kappa\tau)x_k + x_{k+1} = 0$ has general
  solution $x_k = A\sinh(\kappa k\tau) + B\cosh(\kappa k\tau)$,
  fixed by two boundary conditions (same $y_k = r^k$ machinery
  as solving $y_{k-1} - 2 y_k + y_{k+1} = 0$).
  \item \textbf{Tridiagonal linear algebra}: the first-order
  conditions form a tridiagonal system, the matrix version of
  the recurrence.
  \item Mean-variance thinking: minimising $\E[C] + \lambda\Var[C]$
  traces an efficient frontier as $\lambda$ ranges over
  $[0,\infty)$. No schedule is ``optimal'' without specifying
  $\lambda$.
\end{itemize}
\end{backgroundbox}

\section{The liquidation problem}
\label{sec:ac_setup}

\emph{Symbol anchor.} $X$ initial position, $T$ deadline, $N$
intervals, $\tau = T/N$, $\volat$ vol, $\eta$ temporary-impact
slope, $\riskav$ permanent-impact slope, $\lambda$ risk aversion,
$\kappa = \volat\sqrt{\lambda/\eta}$ time-constant.

A \emph{liquidator} holds $X$ shares at $t = 0$ and must hold $0$
at $t = T$, trading at $N+1$ equally spaced times $t_0 = 0 < t_1 <
\ldots < t_N = T$ with spacing $\tau \defeq T/N$. Write $x_k$ for
the \textbf{holding after trade $k$}: $x_0 = X$, $x_N = 0$. The
\textbf{trade at step $k$} is the signed change in holding,
\begin{equation}
  n_k \;\defeq\; x_{k-1} - x_k, \qquad k = 1, 2, \ldots, N.
  \label{eq:ac_nk_def}
\end{equation}
Positive $n_k$ is a sell (holding falls), negative a buy. Trades
sum to the full liquidation,
\begin{equation}
  \sum_{k=1}^{N} n_k \;=\; x_0 - x_N \;=\; X,
  \label{eq:ac_ntotal}
\end{equation}
so the constraint $\sum n_k = X$ is automatic given
$x_0 = X, x_N = 0$; no separate Lagrange multiplier is needed.

\subsection{Three cost sources}
\label{subsec:ac_costs}

Each share sold incurs cost from three distinct mechanisms.
\citeAC{} writes each as an affine function of trading rate $v_k
\defeq n_k/\tau$ (shares per unit time in period $k$), with
separate slopes.

\paragraph{Temporary impact $g(v)$.} Selling $n_k$ shares in the
$k$-th interval consumes visible liquidity on the bid side. The
market maker fills the flow below the no-pressure mid; once the
interval ends the book refills and the temporary component
vanishes. \citeAC{} models this as
\begin{equation}
  g(v) \;=\; \epsilon\,\mathrm{sign}(v) \;+\; \eta\, v,
  \label{eq:ac_temp}
\end{equation}
with two constants:
\begin{itemize}
  \item $\epsilon > 0$: fixed half-spread per share. Every market
  order crosses the bid--ask by at least $\epsilon$. For a pure
  liquidation $\epsilon\sum|n_k| = \epsilon X$ is a constant
  offset and drops out of every optimisation.
  \item $\eta > 0$: linear slope on rate. Twice the rate, twice
  the temporary impact per share. This is the term that makes
  ``rush'' expensive and shapes the optimal schedule.
\end{itemize}
With constant rate over the interval, the average execution price
in period $k$ is the starting mid minus $g(n_k/\tau)$. Total
temporary-impact cost over the session:
\begin{equation}
  C_{\text{temp}} \;=\; \sum_{k=1}^N n_k \cdot g(n_k/\tau)
   \;=\; \epsilon X \;+\; \frac{\eta}{\tau}\sum_{k=1}^N n_k^2.
   \label{eq:ac_ctemp}
\end{equation}
The constant $\epsilon X$ will be dropped everywhere below.

\paragraph{Permanent impact $h(v)$.} Selling signals that somebody
thinks fair value has moved; the market updates its belief and
the mid drifts down. Unlike the temporary component, this drift
does not revert; it lives on the efficient price. \citeAC{}
models it linearly:
\begin{equation}
  h(v) \;=\; \riskav\, v,
  \label{eq:ac_perm}
\end{equation}
with $\riskav > 0$. Over interval $k$ the mid drifts down by
$\tau\cdot h(n_k/\tau) = \riskav n_k$.

\begin{pitfallbox}[A--C $\riskav$ is \emph{not} A--S $\riskav$]
In \cref{ch:market_making}, Avellaneda--Stoikov's reservation price
$r = \Spx - \invq\,\riskav\,\volat^2(T-t)$ used $\riskav$ for the
market-maker's exponential-utility risk aversion. Same Greek, but
here $\riskav$ is the \emph{permanent-impact slope} $h(v) = \riskav
v$, a venue/asset property rather than a trader preference. For the
A--C liquidator, ``risk aversion'' is $\lambda$, introduced below.
Henceforth: $\riskav$ means permanent-impact coefficient, $\lambda$
means liquidator urgency.
\end{pitfallbox}

\paragraph{Price risk $\sigma\sqrt{\tau}\xi_k$.} The mid also moves
for reasons unrelated to the liquidator: other flow, news, macro.
\citeAC{} models this as arithmetic Brownian motion with vol
$\volat$ per unit time, discretised into i.i.d.\ Gaussian shocks:
\begin{equation}
  S_k \;=\; S_{k-1} \;+\; \volat\sqrt{\tau}\,\xi_k
    \;-\; \tau\, h(n_k/\tau),
  \qquad \xi_k \sim \Normal(0,1)\ \iid
  \label{eq:ac_Srecur}
\end{equation}
The Gaussian $\xi_k$ is the random part; $-\tau h(n_k/\tau)$ is
the deterministic permanent-impact drift per interval. \citeAC{}
uses arithmetic rather than geometric Brownian motion because
over a single day the difference is second-order in total cost,
and arithmetic BM admits closed forms where geometric does not.

\begin{keyfactbox}[The six A--C parameters plus one]
Six model parameters plus one decision parameter:
\begin{center}\begin{tabular}{ll}\toprule
  \textbf{Symbol} & \textbf{Meaning} \\ \midrule
  $X$      & Total shares to liquidate \\
  $T$      & Liquidation deadline \\
  $N$      & Number of equal-length trade intervals \\
  $\volat$ & Arithmetic volatility of the underlying per unit time \\
  $\eta$   & Temporary-impact slope (linear part of $g$) \\
  $\riskav$& Permanent-impact slope ($h(v) = \riskav v$) \\ \midrule
  $\lambda$& Liquidator risk aversion (mean-variance penalty) \\
  \bottomrule
\end{tabular}\end{center}
The half-spread $\epsilon$ is a seventh parameter, but since
$\epsilon\sum|n_k| = \epsilon X$ is constant over schedules, it
does not affect the optimal trajectory.
\end{keyfactbox}

\section{The mean-variance objective}
\label{sec:ac_objective}

Total cost of a schedule $\{n_k\}$ (equivalently trajectory
$\{x_k\}$) is a random variable; we form a scalar objective that
trades its mean against its variance. The tension is the whole
point: trading fast pays more impact cost (each child order eats
deeper into the book), trading slow leaves unsold shares exposed
to adverse random drift for longer. Expected cost penalises the
former, variance penalises the latter; A--C picks the schedule
that minimises their weighted sum. The next two subsections unpack
each in turn.

\subsection{Expected cost}
\label{subsec:ac_expectation}

At the start of period $k$ the mid is $S_{k-1}$. The liquidator
sells $n_k$ at average price $S_{k-1} - \eta n_k/\tau$ (dropping
the $\epsilon$ constant), so cash received is $n_k(S_{k-1} - \eta
n_k/\tau)$. The \textbf{implementation shortfall} (``slippage'') is
what the liquidator would have received selling $X$ instantly at
$S_0$ minus what they actually receive:
\begin{equation}
  C \;\defeq\; X\,S_0 \;-\; \sum_{k=1}^N n_k\left(S_{k-1} - \frac{\eta n_k}{\tau}\right).
  \label{eq:ac_cost_def}
\end{equation}
Iterating~\eqref{eq:ac_Srecur} gives $S_{k-1} = S_0 +
\volat\sqrt\tau\sum_{j=1}^{k-1}\xi_j - \riskav\sum_{j=1}^{k-1}n_j$;
expanding the sum yields (\citeAC{} eq.~(4))
\begin{equation}
  C \;=\; \underbrace{\sum_{k=1}^N \tau\,\riskav\,x_k\,\frac{n_k}{\tau}}_{\text{permanent impact}}
       \;+\; \underbrace{\frac{\eta}{\tau}\sum_{k=1}^N n_k^2}_{\text{temporary impact}}
       \;-\; \underbrace{\volat\sqrt{\tau}\sum_{k=1}^N x_k\,\xi_k}_{\text{random price risk}}.
  \label{eq:ac_cost_decomposed}
\end{equation}
Permanent impact is deterministic and quadratic in the trajectory;
temporary impact is deterministic and quadratic in the rates; the
price-risk term is the only random one, linear in Gaussian shocks.

Taking the expectation, only the first two survive:
\begin{equation}
  \E[C] \;=\; \riskav \sum_{k=1}^N x_k\,n_k
         \;+\; \frac{\eta}{\tau}\sum_{k=1}^N n_k^2.
  \label{eq:ac_EC_raw}
\end{equation}

The permanent-impact sum $\sum x_k n_k$ simplifies by a
telescoping identity central to everything below. With
$n_k = x_{k-1} - x_k$,
\begin{equation}
  \sum_{k=1}^N x_k n_k
    \;=\; \sum_{k=1}^N x_k(x_{k-1} - x_k)
    \;=\; \sum_{k=1}^N \left(\tfrac{1}{2}(x_{k-1}^2 - x_k^2) - \tfrac{1}{2}n_k^2\right)
    \;=\; \tfrac{1}{2}(X^2 - 0) - \tfrac{1}{2}\sum_{k=1}^N n_k^2,
  \label{eq:ac_telescope}
\end{equation}
using $x_k(x_{k-1}-x_k) = \tfrac{1}{2}(x_{k-1}^2 - x_k^2) -
\tfrac{1}{2}(x_{k-1}-x_k)^2$ term-by-term and telescoping. Expected
cost becomes purely quadratic in the $n_k$, plus the constant
$\tfrac{1}{2}\riskav X^2$:
\begin{equation}
  \E[C] \;=\; \tfrac{1}{2}\riskav X^2
            \;+\; \frac{1}{\tau}\bigl(\eta - \tfrac{1}{2}\riskav\tau\bigr)
                    \sum_{k=1}^N n_k^2.
  \label{eq:ac_EC_clean}
\end{equation}
The \emph{adjusted temporary-impact coefficient} is
\begin{equation}
  \tilde\eta \;\defeq\; \eta \;-\; \tfrac{1}{2}\riskav\tau,
  \label{eq:ac_etatilde}
\end{equation}
which replaces $\eta$ in the discrete problem. The
$\tfrac{1}{2}\riskav\tau$ correction is the most-missed detail in
re-derivations of A--C: discrete uses $\tilde\eta$, not $\eta$. In
the $\tau \to 0$ limit (\cref{sec:ac_ct_limit}) $\tilde\eta \to
\eta$, so continuous-time uses the un-adjusted $\eta$.

\subsection{Variance of cost}
\label{subsec:ac_variance}

The only random term in~\eqref{eq:ac_cost_decomposed} is
$-\volat\sqrt\tau\sum_k x_k \xi_k$. Its variance:
\begin{equation}
  \Var[C]
  \;=\; \volat^2\tau\,\Var\!\left[\sum_{k=1}^N x_k\xi_k\right]
  \;=\; \volat^2\tau \sum_{k=1}^N x_k^2,
  \label{eq:ac_var}
\end{equation}
by independence of the $\xi_k$. The sum $\sum x_k^2$ is the
discrete $L^2$-norm of the inventory path, so \textbf{variance is
the area under the squared-inventory curve}, weighted by
$\volat^2\tau$. A slow schedule (large $x_k$ throughout) carries
more variance: price moves act on more shares for longer.

\subsection{The objective}
\label{subsec:ac_obj}

\citeAC{} combines them into a linear \textbf{mean-variance
objective}:
\begin{equation}
  U(\{n_k\}) \;\defeq\; \E[C] + \lambda\Var[C],
  \qquad \lambda \ge 0.
  \label{eq:ac_obj}
\end{equation}
Minimise $U$ subject to $x_0 = X, x_N = 0$. $\lambda$ encodes
how much variance the trader accepts to save a unit of expected
cost: large $\lambda$ = urgent, risk-averse (``get me out, I don't
want it overnight''); small $\lambda$ = patient, risk-tolerant
(``minimise cost, I'll wear the mark-to-market risk'').

\begin{keyfactbox}[The A--C objective]
Given $X$ shares over $N$ intervals of length $\tau$, vol $\volat$,
temporary slope $\eta$, permanent slope $\riskav$, and risk
aversion $\lambda \ge 0$, the A--C objective is
\[
  \min_{\{n_k\}}\quad
    \underbrace{\tfrac{1}{2}\riskav X^2 + \frac{\tilde\eta}{\tau}\sum_{k=1}^N n_k^2}_{\E[C]}
    \;+\;
    \lambda\,
    \underbrace{\volat^2\tau\sum_{k=1}^N x_k^2}_{\Var[C]},
\]
with $\tilde\eta = \eta - \tfrac{1}{2}\riskav\tau$,
$n_k = x_{k-1} - x_k$, $x_0 = X$, $x_N = 0$.
\end{keyfactbox}

\begin{intuitionbox}
Two limiting cases before we derive the minimiser.

\textit{$\lambda \to 0$: only expected cost.} The objective is
$\tfrac{\tilde\eta}{\tau}\sum n_k^2$. By Cauchy--Schwarz, $\sum
n_k^2$ at fixed $\sum n_k = X$ is minimised by equal $n_k = X/N$:
\emph{TWAP} (Time-Weighted Average Price), linear inventory.

\textit{$\lambda \to \infty$: only variance.} Minimising
$\volat^2\tau\sum x_k^2$ under $x_0=X, x_N=0$ sends everything at
$k=1$: all $x_k = 0$ thereafter, variance collapses, cost blows up.

Intermediate $\lambda$: deterministic, monotone-decreasing,
\emph{front-loaded} relative to TWAP, trading some expected
cost for variance reduction by shedding the large early positions
that dominate $\sum x_k^2$.
\end{intuitionbox}

\section{Deriving the discrete optimal trajectory}
\label{sec:ac_derivation}

With the objective quadratic in the trajectory, take $\partial/
\partial x_j$ at each interior $j \in \{1, \ldots, N-1\}$, set to
zero, and recognise the finite-difference equation as a
second-order linear recurrence.

\subsection{First-order conditions}
\label{subsec:ac_foc}

Write $U = \tfrac{1}{2}\riskav X^2 + \tfrac{\tilde\eta}{\tau}\sum
n_k^2 + \lambda\volat^2\tau\sum x_k^2$. For interior $x_j$, only
$n_j = x_{j-1} - x_j$ and $n_{j+1} = x_j - x_{j+1}$ depend on it:
\begin{align*}
  \frac{\partial n_j^2}{\partial x_j}
    &= 2\,n_j\,(-1) = -2(x_{j-1} - x_j), \\
  \frac{\partial n_{j+1}^2}{\partial x_j}
    &= 2\,n_{j+1}\,(+1) = 2(x_j - x_{j+1}).
\end{align*}
Adding them,
\begin{equation*}
  \frac{\partial}{\partial x_j}\sum_{k=1}^N n_k^2
  \;=\; -2(x_{j-1} - x_j) + 2(x_j - x_{j+1})
  \;=\; -2(x_{j-1} - 2x_j + x_{j+1}).
\end{equation*}
The variance sum contributes $2\lambda\volat^2\tau\,x_j$; the
$\tfrac{1}{2}\riskav X^2$ constant drops. Hence:
\begin{equation*}
  \frac{\partial U}{\partial x_j}
  \;=\; -\frac{2\tilde\eta}{\tau}(x_{j-1} - 2x_j + x_{j+1})
       \;+\; 2\lambda\volat^2\tau\,x_j
  \;=\; 0.
\end{equation*}
Dividing by $-2\tilde\eta/\tau$ and rearranging,
\begin{equation}
  x_{j-1} \;-\; 2\left(1 + \frac{\lambda\volat^2\tau^2}{2\tilde\eta}\right) x_j
          \;+\; x_{j+1} \;=\; 0,
  \qquad j = 1, 2, \ldots, N-1.
  \label{eq:ac_foc_recur}
\end{equation}
A second-order linear recurrence with constant coefficient. Compare
to the hyperbolic identity
\begin{equation*}
  2\cosh(\kappa\tau) \;=\; e^{\kappa\tau} + e^{-\kappa\tau},
\end{equation*}
and define $\kappa > 0$ by the implicit equation
\begin{equation}
  \boxed{\;
  \cosh(\kappa\tau) \;=\; 1 \;+\; \frac{\lambda\volat^2\tau^2}{2\tilde\eta}
  \;}
  \label{eq:ac_kappa_exact}
\end{equation}
so that~\eqref{eq:ac_foc_recur} becomes, exactly,
\begin{equation}
  x_{j-1} \;-\; 2\cosh(\kappa\tau)\,x_j \;+\; x_{j+1} \;=\; 0.
  \label{eq:ac_recur_sinh}
\end{equation}
Expanding $\cosh(\kappa\tau) = 1 + \tfrac{1}{2}(\kappa\tau)^2 +
O(\tau^4)$ and matching~\eqref{eq:ac_foc_recur} gives the
small-$\tau$ approximation:
\begin{equation}
  \boxed{\;\kappa^2 \;\approx\; \frac{\lambda\volat^2}{\tilde\eta}\;}
  \label{eq:ac_kappa_approx}
\end{equation}
This is the $\kappa$ quoted in textbooks and code. The exact
discrete $\kappa$ from~\eqref{eq:ac_kappa_exact} differs by
$O(\tau^2)$, numerically negligible in practice, and matches the
un-adjusted continuous-time definition as $\tau \to 0$.

\subsection{Solving the recurrence}
\label{subsec:ac_solve_recur}

The recurrence~\eqref{eq:ac_recur_sinh} has characteristic equation
$r^2 - 2\cosh(\kappa\tau) r + 1 = 0$ with roots $r_{\pm} =
e^{\pm\kappa\tau}$, so
\begin{equation*}
  x_k \;=\; A\,e^{k\kappa\tau} + B\,e^{-k\kappa\tau}
       \;=\; A'\,\sinh(\kappa k\tau) + B'\,\cosh(\kappa k\tau).
\end{equation*}
Boundary conditions pin the constants. $x_N = 0$ at $t_N = T$ gives
\begin{equation*}
  0 \;=\; A'\sinh(\kappa T) + B'\cosh(\kappa T)
  \;\Rightarrow\;
  B' \;=\; -A'\tanh(\kappa T).
\end{equation*}
Substituting back,
\begin{align*}
  x_k &= A'\bigl[\sinh(\kappa k\tau) - \tanh(\kappa T)\cosh(\kappa k\tau)\bigr] \\
      &= \frac{A'}{\cosh(\kappa T)}\bigl[\sinh(\kappa k\tau)\cosh(\kappa T) - \sinh(\kappa T)\cosh(\kappa k\tau)\bigr].
\end{align*}
The bracketed expression is a hyperbolic-sine subtraction formula:
$\sinh(\alpha)\cosh(\beta) - \sinh(\beta)\cosh(\alpha) = \sinh(\alpha -
\beta)$, so
\begin{equation*}
  x_k = \frac{A'}{\cosh(\kappa T)}\sinh(\kappa k\tau - \kappa T)
       = -\frac{A'}{\cosh(\kappa T)}\sinh(\kappa(T - k\tau)).
\end{equation*}
Finally, $x_0 = X$ sets $X = -A'\sinh(\kappa T)/\cosh(\kappa T)$, so
$A' = -X\cosh(\kappa T)/\sinh(\kappa T)$, and
\begin{equation}
  \boxed{\;
  x_k \;=\; X \cdot \frac{\sinh\!\bigl(\kappa(T - k\tau)\bigr)}{\sinh(\kappa T)},
  \qquad k = 0, 1, \ldots, N.
  \;}
  \label{eq:ac_xstar}
\end{equation}

\subsection{Trade sizes}
\label{subsec:ac_nstar}

Differencing the trajectory,
\begin{equation*}
  n_k \;=\; x_{k-1} - x_k
       \;=\; \frac{X}{\sinh(\kappa T)}\Bigl[\sinh(\kappa(T - (k-1)\tau))
                             - \sinh(\kappa(T - k\tau))\Bigr].
\end{equation*}
Using $\sinh\alpha - \sinh\beta = 2\cosh\!\bigl(\tfrac{\alpha+\beta}{2}\bigr)
\sinh\!\bigl(\tfrac{\alpha-\beta}{2}\bigr)$ with
$\alpha = \kappa(T - (k-1)\tau)$, $\beta = \kappa(T - k\tau)$, so
$\tfrac{\alpha+\beta}{2} = \kappa(T - (k-\tfrac{1}{2})\tau)$ and
$\tfrac{\alpha-\beta}{2} = \tfrac{1}{2}\kappa\tau$, we get
\begin{equation}
  \boxed{\;
  n_k \;=\; \frac{2\,X\,\sinh\!\bigl(\tfrac{1}{2}\kappa\tau\bigr)}{\sinh(\kappa T)}
       \cdot \cosh\!\bigl(\kappa\bigl(T - (k-\tfrac{1}{2})\tau\bigr)\bigr),
  \qquad k = 1, 2, \ldots, N.
  \;}
  \label{eq:ac_nstar}
\end{equation}
The factor $\cosh(\kappa(T - (k-\tfrac{1}{2})\tau))$ decreases
monotonically from $\cosh(\kappa(T-\tau/2))$ at $k=1$ to
$\cosh(\kappa\tau/2)$ at $k=N$: trades shrink over time, the
schedule front-loads. The prefactor $2X\sinh(\kappa\tau/2)/
\sinh(\kappa T)$ normalises the total to $X$ by telescoping.

\begin{keyfactbox}[Almgren--Chriss optimal trajectory]
The inventory and trade-size schedules that minimise the
mean-variance objective $\E[C] + \lambda\Var[C]$ subject to
$x_0 = X$, $x_N = 0$ are
\[
  x_k \;=\; X\,\frac{\sinh\!\bigl(\kappa(T - k\tau)\bigr)}{\sinh(\kappa T)},
  \qquad
  n_k \;=\; \frac{2X\sinh(\tfrac{1}{2}\kappa\tau)}{\sinh(\kappa T)}\,\cosh\!\bigl(\kappa(T - (k-\tfrac{1}{2})\tau)\bigr),
\]
with
\[
  \kappa^2 \;=\; \frac{\lambda\volat^2}{\tilde\eta},
  \qquad
  \tilde\eta \;=\; \eta - \tfrac{1}{2}\riskav\tau.
\]
Boundary conditions are automatic: $x_0 = X$ because
$\sinh(\kappa T)/\sinh(\kappa T) = 1$, and $x_N = 0$ because
$\sinh(0) = 0$.
\end{keyfactbox}

\section{Continuous-time limit}
\label{sec:ac_ct_limit}

Take $\tau \to 0, N \to \infty$ with $N\tau = T$ fixed:
\begin{itemize}
  \item $\tilde\eta = \eta - \tfrac{1}{2}\riskav\tau \to \eta$.
  \item $\kappa^2 = \lambda\volat^2/\tilde\eta \to
  \lambda\volat^2/\eta$, a function of continuous parameters only.
  \item $x_k \to x(t)$, continuous in $t \in [0, T]$.
\end{itemize}
The discrete solution~\eqref{eq:ac_xstar} becomes
\begin{equation}
  \boxed{\;
  x(t) \;=\; X\,\frac{\sinh\!\bigl(\kappa(T - t)\bigr)}{\sinh(\kappa T)},
  \qquad \kappa \;=\; \volat\sqrt{\frac{\lambda}{\eta}}.
  \;}
  \label{eq:ac_ct}
\end{equation}
\citeAC{}~\S4 derives the continuous problem directly via
Euler--Lagrange on $\int_0^T[\eta\dot{x}^2 + \lambda\volat^2
x^2]\,\mathrm{d}t$, giving $\eta\ddot{x} - \lambda\volat^2 x = 0$
and hence~\eqref{eq:ac_ct}; the two derivations agree.

\section{Worked numerical example}
\label{sec:ac_worked}

A plug-in is worth a page of formulas. We take the canonical
mid-cap US equity liquidation that appears in every textbook.

\paragraph{Parameters.}
\begin{itemize}
  \item $X = 10^6$ shares (a round lot).
  \item $T = 1$ trading day, $N = 13$ half-hour intervals,
  $\tau \approx 0.0769$ day.
  \item $\volat = 0.30$ fractional-price per $\sqrt{\text{day}}$;
  for a $\$50$ stock, daily stdev $\approx \$15$.
  \item $\eta = 2.5\times 10^{-6}$. At rate $10^7$ shares/day the
  concession is $\eta v = 25$ (huge, because rate is huge); at a
  realistic $X/T = 10^6$/day it is $2.5$ price units, a plausible
  heavy-liquidation magnitude.
  \item $\riskav = 2.5\times 10^{-7}$. Permanent impact over the
  whole day is $\riskav X = 0.25$, an order of magnitude below
  the temporary cost, matching the empirical observation that
  permanent impact is a minority of execution cost.
  \item $\lambda = 10^{-6}$, mildly risk-averse.
\end{itemize}

\paragraph{Derived quantities.} From~\eqref{eq:ac_etatilde},
\begin{equation*}
  \tilde\eta \;=\; \eta - \tfrac{1}{2}\riskav\tau
              \;=\; 2.5\times 10^{-6} - \tfrac{1}{2}(2.5\times 10^{-7})(0.0769)
              \;=\; 2.4904\times 10^{-6}.
\end{equation*}
From~\eqref{eq:ac_kappa_approx},
\begin{equation*}
  \kappa^2 \;=\; \frac{\lambda\volat^2}{\tilde\eta}
           \;=\; \frac{10^{-6}\cdot 0.09}{2.4904\times 10^{-6}}
           \;=\; 0.03614,
  \qquad
  \kappa \;=\; 0.1901.
\end{equation*}
With $T = 1$, $\kappa T = 0.1901$ and $\sinh(\kappa T) = 0.19125$.

\paragraph{Trajectory.} Applying~\eqref{eq:ac_xstar},

\begin{center}
\begin{tabular}{rrrr}\toprule
  $k$ & $t_k$ (day) & $x_k$ (shares) & $n_k$ (shares) \\ \midrule
  $0$ & $0.0000$ & $1\,000\,000.00$ & --- \\
  $1$ & $0.0769$ & $922\,256.65$ & $77\,743.35$ \\
  $2$ & $0.1538$ & $844\,710.53$ & $77\,546.13$ \\
  $3$ & $0.2308$ & $767\,345.04$ & $77\,365.49$ \\
  $4$ & $0.3077$ & $690\,143.64$ & $77\,201.40$ \\
  $5$ & $0.3846$ & $613\,089.83$ & $77\,053.81$ \\
  $6$ & $0.4615$ & $536\,167.12$ & $76\,922.71$ \\
  $7$ & $0.5385$ & $459\,359.06$ & $76\,808.05$ \\
  $8$ & $0.6154$ & $382\,649.24$ & $76\,709.82$ \\
  $9$ & $0.6923$ & $306\,021.25$ & $76\,627.99$ \\
  $10$ & $0.7692$ & $229\,458.69$ & $76\,562.55$ \\
  $11$ & $0.8462$ & $152\,945.21$ & $76\,513.49$ \\
  $12$ & $0.9231$ & $76\,464.43$ & $76\,480.78$ \\
  $13$ & $1.0000$ & $0.00$ & $76\,464.43$ \\ \bottomrule
\end{tabular}
\end{center}

\textbf{Boundary checks.} $x_0 = 10^6$ and $x_{13} = 0.00$ to the
precision shown; $\sum_{k=1}^{13} n_k = 1\,000\,000.00 = X$ by
construction.

\textbf{Vs.\ TWAP.} TWAP sends $X/N \approx 76\,923.08$ per bin;
A--C's first trade is $77\,743.35$ ($+1.07\%$ vs.\ TWAP), its
last $76\,464.43$ ($-0.60\%$). At $\lambda = 10^{-6}$, A--C is
almost indistinguishable from TWAP: $\kappa T = 0.19$ is small,
so $\sinh(\kappa(T-t))/\sinh(\kappa T) \approx (T-t)/T$. Curvature
emerges at higher $\lambda$ (see \cref{sec:ac_frontier}). For a
gently risk-averse trader, A--C says ``don't deviate from TWAP''.

\begin{intuitionbox}
The smallness of $\kappa T$ is not a parameter accident: for
one-day liquidations of liquid names, linear impact and vol balance
so $\kappa T \ll 1$, and optimal execution is only mildly
front-loaded vs.\ TWAP. A--C's main payoff over TWAP is not a
dramatic schedule change but a principled way to pick
$\lambda$ from client risk tolerance, and to compute the chosen
schedule's expected cost and standard deviation for post-trade
benchmarking.
\end{intuitionbox}

\section{The execution efficient frontier}
\label{sec:ac_frontier}

Sweeping $\lambda \in [0,\infty)$ at fixed $X, T, N, \volat,
\eta, \riskav$ traces a curve in the $(\sqrt{\Var[C]}, \E[C])$
plane: the \textbf{execution efficient frontier}, a Markowitz-style
curve of minimum cost at given variance (or vice versa). The analogy
is exact: where Markowitz portfolio selection trades expected return
against variance over asset weights, A--C trades expected cost
against variance over execution schedules. Each point on the curve
is the best achievable $\E[C]$ at its $\sqrt{\Var[C]}$; no schedule
lies below it, and picking $\lambda$ just picks which point on the
curve you want.

\begin{figure}[ht]
\centering
\begin{tikzpicture}
\begin{axis}[
  width=10cm, height=6.5cm,
  xlabel={Variance of execution cost $\mathrm{Var}[\mathrm{cost}]$},
  ylabel={Expected execution cost $\mathbb{E}[\mathrm{cost}]$},
  xmin=0, xmax=1.1,
  ymin=0, ymax=1.1,
  axis lines=left,
  major tick length=0pt,
  xtick={0,0.5,1},
  xticklabels={0, medium risk, high risk},
  ytick={0,0.5,1},
  yticklabels={0, , high cost},
  title={Almgren--Chriss efficient frontier},
  title style={font=\small\bfseries},
]
\addplot[blue!70, thick] coordinates {
  (0.04,1.0)(0.08,0.80)(0.15,0.58)};
\addplot[blue!70, thick] coordinates {
  (0.15,0.58)(0.25,0.42)(0.38,0.31)};
\addplot[blue!70, thick] coordinates {
  (0.38,0.31)(0.53,0.24)(0.70,0.20)(0.88,0.17)(1.05,0.16)};
\node[circle, fill=red!70, inner sep=2.5pt]
  at (axis cs:0.53,0.24) {};
\node[font=\footnotesize, above right=2pt, red!80]
  at (axis cs:0.53,0.24) {TWAP};
\node[circle, fill=green!60!black, inner sep=2.5pt]
  at (axis cs:0.08,0.80) {};
\node[font=\footnotesize, above right=2pt, green!60!black, align=center]
  at (axis cs:0.08,0.80) {fast execution\\(high impact)};
\node[circle, fill=orange!80, inner sep=2.5pt]
  at (axis cs:0.88,0.17) {};
\node[font=\footnotesize, above=5pt, orange!80, align=center]
  at (axis cs:0.88,0.17) {slow execution\\(high price risk)};
\draw[-Stealth, dashed, gray!70]
  (axis cs:0.62,0.56) -- (axis cs:0.18,0.70)
  node[midway, above, sloped, font=\footnotesize]{$\lambda\uparrow$};
\end{axis}
\end{tikzpicture}
\caption{Almgren--Chriss efficient frontier. Each point on the curve
corresponds to a different risk-aversion parameter $\lambda$: high $\lambda$
forces fast execution (low variance from price risk, high expected cost
from market impact); low $\lambda$ allows slow execution (low impact cost,
but high variance as price may drift). TWAP sits near the middle of the
frontier. The optimal schedule for a given $\lambda$ is the point on this
curve that minimises $\mathbb{E}[\mathrm{cost}] + \lambda\,\mathrm{Var}[\mathrm{cost}]$.}
\label{fig:ac_frontier}
\end{figure}

Endpoints: $\lambda = 0$ sits at the rightmost (highest-variance,
lowest-cost) point, TWAP; $\lambda \to \infty$ sits at the
leftmost, the instant block trade. Interior points are
hyperbolic-sine schedules.

For the worked example parameters, using
\begin{align*}
  \E[C]   &\;=\; \tfrac{1}{2}\riskav X^2 + \frac{\tilde\eta}{\tau}\sum_k n_k^2, \\
  \Var[C] &\;=\; \volat^2\tau\sum_k x_k^2,
\end{align*}
we compute:

\begin{center}
\begin{tabular}{rrrr}\toprule
  $\lambda$ & $\E[C]$ (units) & $\sqrt{\Var[C]}$ (units) & $n_1$ as \% of $X$ \\ \midrule
  $10^{-7}$ & $2{,}615{,}385$ & $163{,}135$ & $7.70\%$ \\
  $10^{-6}$ & $2{,}615{,}456$ & $162{,}740$ & $7.77\%$ \\
  $10^{-5}$ & $2{,}622{,}090$ & $158{,}944$ & $8.49\%$ \\
  $10^{-4}$ & $3{,}017{,}007$ & $131{,}389$ & $14.28\%$ \\
  $10^{-3}$ & $7{,}481{,}151$ & $67{,}447$  & $37.03\%$ \\ \bottomrule
\end{tabular}
\end{center}

Increasing $\lambda$ drives $\sqrt{\Var[C]}$ down (good) and
$\E[C]$ up (bad). At $\lambda = 10^{-3}$ the first trade is over
a third of the position: the trader pays nearly $3\times$ the
expected cost to shed overnight variance. The desk's job is to
pick the point matching client risk tolerance; A--C maps the
qualitative ``client wants out fast'' to a specific $\lambda$ and
hence a specific schedule.

\begin{intuitionbox}
Two observations. First, the permanent-impact term
$\tfrac{1}{2}\riskav X^2 = 1.25\times 10^5$ is small vs.\ the
$\sim 2.5\times 10^6$ expected-cost totals; most cost is temporary
impact from $\sum n_k^2$. In this regime the $\tilde\eta$-vs-$\eta$
distinction is a small correction, but it is the one that makes the
discrete closed form exact.

Second, at $\lambda = 10^{-3}$ the trader pays an extra $\$4.9$m
expected cost to save $\$95$k of cost standard deviation.
Extreme risk aversion; no desk would pick it for a routine
liquidation, but it fits regimes like an executing broker with a
hard deadline and a client indifferent to average cost. The
frontier is a menu, not a prescription.
\end{intuitionbox}

\section{VWAP, TWAP, POV, and A--C in context}
\label{sec:ac_taxonomy}

Sell-side execution algorithms come in a small zoo. Four matter
for this book.

\paragraph{TWAP (Time-Weighted Average Price).} Equal size
per interval, $n_k = X/N$. The name: when price is martingale, fill
VWAP equals the mid's TWAP. TWAP is the $\lambda \to 0$ limit of
A--C, default for ``lowest expected cost, don't care about timing
risk''. Simple, benchmark-friendly, gaming-proof under martingale
mid.

\paragraph{VWAP (Volume-Weighted Average Price).} Trade in
proportion to the \emph{historical} intraday volume profile ($10\%$
of daily volume in the first half-hour $\Rightarrow 10\%$ of the
order). Benchmark: session realised VWAP. The default institutional
equity algorithm because it mimics a ``typical'' participant and
minimises tracking error to the market's average print. Not a
direct A--C special case, but A--C re-weighted by a volume profile
reproduces it.

\paragraph{POV (Percentage of Volume).} Trade a fixed fraction
$\rho$ of \emph{realised} volume. If the last minute printed
$10{,}000$ shares at $\rho = 5\%$, send $500$ next minute. Reactive:
speeds up with the market, slows with it, so the liquidator's
footprint stays proportional. Default for hiding in flow at the
cost of deadline risk (insufficient volume $\Rightarrow$ incomplete
fill).

\paragraph{A--C / Implementation Shortfall.} Solve this chapter's
trajectory problem for given $\lambda$ and execute the
hyperbolic-sine schedule. The answer to ``minimise implementation
shortfall (starting mid minus average fill) subject to a risk
penalty''. Front-loaded vs.\ TWAP; equivalent to VWAP plus a risk
correction under a flat volume profile. A--C (or a risk-aware
cousin) is every modern desk's ``Implementation Shortfall''
flagship.

\begin{keyfactbox}[Execution algorithm taxonomy]
\begin{center}\footnotesize\begin{tabular}{>{\raggedright\arraybackslash}p{1.3cm}
                                 >{\raggedright\arraybackslash}p{5.8cm}
                                 >{\raggedright\arraybackslash}p{5.8cm}}\toprule
  \textbf{Algo} & \textbf{Schedule rule} & \textbf{Benchmark} \\ \midrule
  TWAP & Equal per time bin & Session mid TWAP \\
  VWAP & Proportional to historical volume profile & Session realised VWAP \\
  POV  & Fixed \% of realised volume live & Realised VWAP (participation-matched) \\
  A--C & Hyperbolic-sine, front-loaded by $\lambda$ & Arrival (decision) mid, i.e.\ implementation shortfall \\
  \bottomrule
\end{tabular}\end{center}
All four are deterministic given inputs; none uses a private
signal. Signal-driven ``algo wheels'' sit on top of this layer.
\end{keyfactbox}

\section{Almgren--Chriss versus square-root reality}
\label{sec:ac_vs_sqrt}

The chapter's central pedagogical point, and the hardest to
swallow first time. A--C models temporary impact linearly,
$g(v) = \eta v$. \Cref{ch:market_impact} showed the empirical
law, stable across decades and asset classes, is square-root:
\begin{equation}
  I(Q) \;\approx\; Y\volat\sqrt{Q/V_{\text{daily}}},
  \label{eq:ch18_sqrt_recall}
\end{equation}
where $Q$ is metaorder size, $V_{\text{daily}}$ ADV, and $Y \sim
0.5$ a dimensionless prefactor. Traded at constant rate over
interval $\tau$, linear $g(v)$ gives $I(Q) = \eta Q/\tau \propto
Q$ vs.\ the square-root's $I(Q) \propto \sqrt{Q}$. The two curves
meet at exactly one order size (where $\eta$ was calibrated)
and diverge elsewhere.

\textbf{Why teach linear anyway?} Three reasons, none being
``linear is empirically right'':
\begin{enumerate}
  \item \textbf{Closed form.} Linear impact yields the quadratic
  objective~\eqref{eq:ac_obj}, linear first-order conditions, a
  hyperbolic-sine solution, and a closed-form conic frontier.
  Square-root impact gives a non-quadratic objective with no
  closed form: a nonlinear BVP to solve numerically. For teaching
  the structure of the cost--risk trade-off, the closed form wins.
  \item \textbf{Local calibration.} In a narrow band around the
  trader's typical size $Q_0$, linear can be tuned to match the
  square-root slope: pick $\eta$ so $\eta Q_0/T = Y\volat\sqrt{Q_0/V}$.
  Within a factor of $2$--$3$ of $Q_0$ it is not catastrophically
  wrong.
  \item \textbf{Industry inertia.} A--C became the industry
  benchmark within five years. Serious desks run non-linear
  engines, but A--C remains the lingua franca and
  $(\eta, \riskav, \lambda)$ are what daily risk reports plot.
\end{enumerate}

\begin{pitfallbox}[Don't extrapolate A--C outside its calibration range]
The cost of this contradiction is largest when a small-trade
calibration is used to price a large one. Suppose $\eta$ is
calibrated against child orders of $Q_{\text{calib}} = 10^5$ shares:
\begin{itemize}
  \item Child $Q = Q_{\text{calib}}$: predicted cost matches reality.
  \item Parent $Q = 5\cdot 10^6$: linear predicts $\propto Q$,
  reality is $\propto \sqrt{Q}$, so linear overpredicts by
  $\sqrt{Q/Q_{\text{calib}}} = \sqrt{50} \approx 7.07$.
\end{itemize}
Running A--C with that over-stated slope tells the trader any
parent liquidation is catastrophic, steering them to artificially
slow schedules and ballooning variance. Fix: recalibrate $\eta$
at the parent-order size band before running A--C, never scale up
a child calibration. Empirical backing: \cref{ch:market_impact}.
\end{pitfallbox}

\section{Extensions and breakages}
\label{sec:ac_extensions}

Directions the basic model generalises in; doorways to further
reading, not rederived here.

\paragraph{Risk-neutral vs.\ risk-averse.} $\lambda = 0$ (TWAP as
A--C solution) is the risk-neutral case; some authors treat it as
a separate ``pure implementation shortfall'' problem. Replacing
mean-variance with a full utility $U(C) = -\exp(\gamma C)$ gives
the same closed form under Gaussian risk (CARA--Gaussian is
mean-variance up to constants); CVaR or VaR break the closed form.

\paragraph{Time-varying volatility.} $\volat \to \volat(t)$
(higher at open/close, lower midday) generalises the trajectory:
$\eta\ddot{x}(t) - \lambda\volat(t)^2 x(t) = 0$, a variable-coefficient
ODE, numerical but not generally elementary. See \citeCartea{}
Ch.~6.

\paragraph{Multi-asset portfolio liquidation.} Vector positions
$\bm{X} \in \R^n$ with covariance $\bm{\Sigma}$ and cross-impact
$\bm{\eta}$ give a matrix quadratic with hyperbolic-sine matrix
$\bm{\kappa} = \bm{\eta}^{-1/2}(\lambda\bm{\Sigma})^{1/2}\bm{\eta}^{-1/2}$.
Closed form survives; bookkeeping doesn't. See \citeCartea{}
Ch.~7.

\paragraph{Non-linear impact.} With $g(v) \propto \sqrt{v}$ the
objective is non-quadratic and has no closed form. Obizhaeva--Wang
(2013) and Almgren (2003) give numerical treatments;
large-order practitioners run non-linear optimisers on proprietary
square-root curves.

\paragraph{Dark pools and venue choice.} Dark (non-displayed)
execution cuts temporary impact at the cost of uncertain fills.
A--C with venue choice is an open research area; the simplest
treatment mixes displayed and dark $\eta$'s with a dark fill
probability. See \citeBouchaud{}.

\section{Prosperity and sell-side applications}
\label{sec:ac_applications}

\begin{prosperitybox}
End-of-round liquidation is where A--C earns its keep in
Prosperity. Suppose you are long $30$ PICNIC\_BASKET1 entering the
final $100$ timesteps after the basket--constituent spread you
were harvesting has collapsed. You want flat by $T = 100$: the
round ends and unliquidated inventory marks to the last mid
regardless of realisable exit. You also want to avoid dumping all
$30$ at once: the basket is thin and a market order of that
size walks the book down several ticks per unit.

Take parameters $X = 30$, $T = 100$, $N = 10$ (one decision every
$10$ timesteps), $\volat = 2.0$ seashells per timestep,
$\eta = 0.05$ (each unit/timestep of selling rate drives the mid
down by $0.05$), $\riskav = 0.001$ (permanent impact is small on
the order book), and a mild $\lambda = 10^{-6}$. The A--C
trajectory is computed exactly as in \cref{sec:ac_worked}:
\[
  \tilde\eta = 0.05 - \tfrac{1}{2}\cdot 0.001\cdot 10 = 0.045,
  \qquad \kappa = \sqrt{10^{-6}\cdot 4 / 0.045} \approx 0.00943,
\]
so $\kappa T \approx 0.943$, larger than the equity example's
$0.19$: the trader is closer to the deadline in P\&L-standard-deviation
units. The schedule:
\begin{center}
\begin{tabular}{rrrr}\toprule
  $k$ & $t_k$ & $x_k$ & $n_k$ \\ \midrule
  $0$  & $0$   & $30.00$ & --- \\
  $1$  & $10$  & $26.29$ & $3.71$ \\
  $2$  & $20$  & $22.81$ & $3.48$ \\
  $3$  & $30$  & $19.53$ & $3.28$ \\
  $4$  & $40$  & $16.43$ & $3.10$ \\
  $5$  & $50$  & $13.48$ & $2.96$ \\
  $6$  & $60$  & $10.64$ & $2.84$ \\
  $7$  & $70$  & $7.90$  & $2.74$ \\
  $8$  & $80$  & $5.23$  & $2.67$ \\
  $9$  & $90$  & $2.60$  & $2.63$ \\
  $10$ & $100$ & $0.00$  & $2.60$ \\ \bottomrule
\end{tabular}
\end{center}
TWAP sends $3.00$ per decision; A--C sends $3.71$ first ($+24\%$)
and $2.60$ last ($-13\%$), trading a little early impact for tail
variance reduction. $\sum n_k = 30.000$ exactly by construction.

\textbf{Prosperity caveats.} (1) Position limits: if the limit is
$\pm 30$, A--C does not help when flipping from $+30$ to $-30$;
that is reversing, not liquidating. (2) Discrete units: round
$n_k$ to integers and hold the slack in $x_k$; do not let rounding
accumulate into a non-zero terminal. (3) Mid vs fill: A--C uses
mid; Prosperity fills cross the spread, so the model slightly
over-estimates cost.
\end{prosperitybox}

\begin{gsbox}
On a Goldman execution desk, the default algorithm for a client
order tagged ``IS'' (Implementation Shortfall) is an A--C
descendant with venue routing, a dark-pool seeker, and a
signal-based accelerator overlay. The owning \emph{Strat} does
two unglamorous jobs. First, daily calibration of $\eta$ and
$\riskav$ against the previous day's fills across tens of
thousands of orders, regressing realised shortfall against $(Q/V)$
and $(Q/T)^\alpha$ curves. Second, maintain $\lambda$ presets
(``urgent'', ``neutral'', ``passive'') with a one-click override
so traders pick at order entry without knowing the math. When
asked ``why did my fill trail VWAP today?'' the Strat pulls a
post-trade report decomposing realised shortfall into four
A--C buckets: permanent impact, temporary impact, timing risk,
spread. These reports are the team's daily bread.

A second use is \emph{pre-trade cost estimation}: before a block,
sales ask Strats ``how much to execute in two hours''. The answer
is A--C $\E[C]$ at chosen $\lambda$, with $\sqrt{\Var[C]}$. It
goes on the ticket the salesperson quotes. Getting it
systematically wrong loses clients and gets the Strat paged.
\end{gsbox}

\section{A short reference implementation}
\label{sec:ac_code}

Twenty lines of Python compute the whole trajectory. It is not
production code (real engines need volume profiles, queue-position
logic, child-order sizing against book depth, venue routing), but it
is the core calculation underneath any A--C execution stack.

\begin{lstlisting}[caption={Almgren--Chriss optimal liquidation trajectory.}]
import math

def almgren_chriss_trajectory(X, T, N, sigma, eta, gamma, lam):
    """Return the list of inventory levels [x_0, x_1, ..., x_N]
    under the Almgren-Chriss optimal mean-variance schedule.

    Parameters
    ----------
    X     : total shares to liquidate (x_0 = X, x_N = 0).
    T     : deadline (same time units as sigma, eta, gamma).
    N     : number of equal-length trade intervals.
    sigma : arithmetic volatility per unit time.
    eta   : temporary-impact slope (linear in trading rate).
    gamma : permanent-impact slope (linear in trading rate).
    lam   : mean-variance risk aversion (>= 0).
    """
    tau = T / N
    eta_tilde = eta - 0.5 * gamma * tau
    kappa = math.sqrt(lam * sigma * sigma / eta_tilde)
    denom = math.sinh(kappa * T)
    return [X * math.sinh(kappa * (T - k * tau)) / denom
            for k in range(N + 1)]
\end{lstlisting}

Recover trade sizes by differencing: $n_k = x_{k-1} - x_k$. The
$\lambda = 0$ case throws (\texttt{kappa = 0}, \texttt{sinh(0)} in
the denominator); production code should trap it and return the
TWAP schedule $[X(1 - k/N)]$ directly.

\begin{historybox}
Robert Almgren (Courant, NYU) and Neil Chriss wrote ``Optimal
Execution of Portfolio Transactions'' in 1999; it appeared in the
\emph{Journal of Risk} in 2000 and sits in every execution desk's
literature folder. It was the first clean statement of the
mean-variance execution problem; the hyperbolic-sine closed form
turned ``how fast do we hit the button'' from craft into
discipline. Almgren later founded Quantitative Brokers. Chriss
had sell- and buy-side roles including Goldman Sachs Asset
Management, Hutchin Hill, and Paloma. Newer models
(Obizhaeva--Wang transient impact, Gatheral's no-dynamic-arbitrage
framework) all frame themselves as departures from A--C.
\end{historybox}

\section{Key facts to memorise}
\label{sec:ac_keyfacts}

\begin{itemize}
  \item Problem: sell $X$ over $[0,T]$ with $\sum n_k = X$,
  minimising $\E[C] + \lambda\Var[C]$.
  \item Three cost sources: temporary $g(v) = \eta v$ (reverts),
  permanent $h(v) = \riskav v$ (sticks), price risk
  $\volat\sqrt\tau\xi_k$ (Gaussian).
  \item Telescoping $\sum x_k n_k = \tfrac{1}{2}(X^2 - \sum n_k^2)$
  makes permanent impact quadratic in $n_k$.
  \item Discrete: $\tilde\eta = \eta - \tfrac{1}{2}\riskav\tau$,
  $\kappa^2 = \lambda\volat^2/\tilde\eta$. Continuous limit:
  $\tilde\eta \to \eta$, $\kappa = \volat\sqrt{\lambda/\eta}$.
  \item Optimal trajectory: $x_k = X\sinh(\kappa(T - k\tau))/
  \sinh(\kappa T)$.
  \item $\lambda \to 0$: TWAP. $\lambda \to \infty$: instant
  block. Interior $\lambda$: front-loaded.
  \item Efficient frontier: $\lambda \mapsto (\sqrt{\Var[C]},
  \E[C])$ traces a convex curve.
  \item Four baseline algos: TWAP, VWAP, POV, A--C. A--C is the
  ``Implementation Shortfall'' default.
  \item Linear impact contradicts the square-root law
  (\cref{ch:market_impact}). Kept for tractability; practitioners
  recalibrate $\eta$ locally to the square-root slope.
  \item Goldman desk: Strat owns $(\eta, \riskav)$ calibration
  against live fills; trader picks $\lambda$ from client risk
  tolerance; overrides on news.
  \item A--C $\riskav$ = permanent-impact slope, not the A--S
  risk aversion of \cref{ch:market_making}.
  \item Primary: \citeAC{}. Richer re-derivation: \citeCartea{}
  Ch.~6--8. Empirical counter: \cref{ch:market_impact},
  \citeBouchaud{}.
\end{itemize}
